\chapter{Results: NL-to-SQL Reproduction}\label{ch:results-sql}

All tables in this chapter are generated from the result files (\fpath{report/tools/build_tables.py}, generated
\tablesgenerated). Scores are repo-faithful execution accuracy unless marked ``corrected''.

\section{Vanilla baselines and the gate}\label{sec:sql-vanilla}

{\footnotesize
\rowcolors{3}{rowshade}{white}
\setlength{\tabcolsep}{4.2pt}
\begin{longtable}{@{}lllrrcrrr@{}}
\caption{NL-to-SQL vanilla baselines on the four rebuilt test sets, both checkpoints and both prompts. Execution accuracy (fraction of cases). The paper's vanilla figures (Table~3, Llama-3.1-8B) are set beside our zero-shot rows only, because the paper does not state its test prompt.}\label{tab:gen-vanilla-sql}\\
\toprule
\rowcolor{headshade}Model & Prompt & Test set & $n$ & Repo-faithful & \makecell{95\%\\Wilson CI} & Corrected & \makecell{Paper\\vanilla} & Ours $-$ paper\\
\midrule
\endfirsthead
\multicolumn{9}{l}{\textit{\tablename~\thetable{} (continued)}}\\
\toprule
\rowcolor{headshade}Model & Prompt & Test set & $n$ & Repo-faithful & \makecell{95\%\\Wilson CI} & Corrected & \makecell{Paper\\vanilla} & Ours $-$ paper\\
\midrule
\endhead
\midrule
\multicolumn{9}{r}{\textit{continued on next page}}\\
\endfoot
\bottomrule
\endlastfoot
Base & zero-shot & WikiSQL & 1000 & 0.3260 & [0.2980, 0.3560] & 0.3260 & 0.525 & -0.1990\\
Base & zero-shot & WikiTQ & 1000 & 0.2190 & [0.1940, 0.2460] & 0.2190 & 0.300 & -0.0810\\
Base & zero-shot & Spider & 1198 & 0.7245 & [0.6990, 0.7490] & 0.7529 & 0.670 & +0.0545\\
Base & zero-shot & BIRD & 363 & 0.3636 & [0.3160, 0.4140] & 0.3636 & 0.388 & -0.0244\\
Base & few-shot & WikiSQL & 1000 & 0.6850 & [0.6560, 0.7130] & 0.6850 &  & --\\
Base & few-shot & WikiTQ & 1000 & 0.3440 & [0.3150, 0.3740] & 0.3440 &  & --\\
Base & few-shot & Spider & 1198 & 0.7688 & [0.7440, 0.7920] & 0.7721 &  & --\\
Base & few-shot & BIRD & 363 & 0.3994 & [0.3500, 0.4510] & 0.3994 &  & --\\
Instruct & zero-shot & WikiSQL & 1000 & 0.5180 & [0.4870, 0.5490] & 0.5180 & 0.525 & -0.0070\\
Instruct & zero-shot & WikiTQ & 1000 & 0.2890 & [0.2620, 0.3180] & 0.2890 & 0.300 & -0.0110\\
Instruct & zero-shot & Spider & 1198 & 0.7913 & [0.7670, 0.8130] & 0.7913 & 0.670 & +0.1213\\
Instruct & zero-shot & BIRD & 363 & 0.4022 & [0.3530, 0.4530] & 0.4050 & 0.388 & +0.0142\\
Instruct & few-shot & WikiSQL & 1000 & 0.7980 & [0.7720, 0.8220] & 0.7980 &  & --\\
Instruct & few-shot & WikiTQ & 1000 & 0.4230 & [0.3930, 0.4540] & 0.4230 &  & --\\
Instruct & few-shot & Spider & 1198 & 0.7963 & [0.7730, 0.8180] & 0.7988 &  & --\\
Instruct & few-shot & BIRD & 363 & 0.4187 & [0.3690, 0.4700] & 0.4187 &  & --\\
Base (original-case tables) & zero-shot & WikiSQL & 1000 & 0.1420 & [0.1220, 0.1650] & 0.1420 &  & \\
Base (original-case tables) & few-shot & WikiSQL & 1000 & 0.2870 & [0.2600, 0.3160] & 0.2870 &  & \\
\end{longtable}
}
\vspace{-0.6\baselineskip}
{\footnotesize\color{slate}\textit{Note.} Sources: \fpath{results/baseline/final_*_base.json}, \fpath{gate_*_base_fewshot.json}, \fpath{gate_instruct_*.json}, \fpath{gate_wikisql_base_origcase_*.json}. Base queried with raw completions (128 tokens); Instruct through the repo's chat call (256 tokens).\par}
\medskip


\begin{observation}
\begin{itemize}
\item \textbf{Base, zero-shot} (\cref{tab:gen-vanilla-sql}) meets the $\pm0.05$ gate only on BIRD ($-0.024$). It misses
WikiSQL by $-0.199$ and WikiTQ by $-0.081$, and overshoots Spider by $+0.055$.
\item \textbf{Base, few-shot} meets it on WikiTQ ($+0.044$) and BIRD ($+0.011$), but overshoots WikiSQL (0.685) and
Spider (0.769).
\item \textbf{Instruct}, through the repository's chat call with the zero-shot prompt, is within 0.015 of the paper on
WikiSQL (0.518), WikiTQ (0.289) and BIRD (0.4022). It overshoots Spider by $+0.121$.
\item \textbf{Instruct, few-shot} is much higher: WikiSQL 0.798, WikiTQ 0.423.
\end{itemize}
\end{observation}

\begin{interpretation}
Three facts point to the paper's vanilla column coming from Instruct through the repository's chat call with the
zero-shot prompt:
\begin{itemize}
\item it is the only combination within the gate on three sets;
\item it is the repository's only query path;
\item the zero-shot prompt is the repository's default test prompt.
\end{itemize}
This cannot be confirmed without the authors. Spider is above the paper for every model and prompt of ours, so the
Spider gap is attributed to the case set or the scoring, not to the model. Our Spider count (1,198) matches the paper
exactly, but its cases cannot be checked.
\end{interpretation}

On base WikiSQL, most failures are execution errors: 507 of 1,000 zero-shot answers fail to run, often because a
multi-word column name is unquoted (\texttt{WHERE Last Meeting = ...}). The few-shot prompt fixes the quoting (25
failures), which explains most of its +0.359. Showing the WikiSQL tables in their original case lowered base to 0.142
(zero-shot) and 0.287 (few-shot), because the database values are lowercase. The paper's WikiSQL score falls between
the two prompts, and no tested combination lands near it.

\section{The core study on base: TS against NV}\label{sec:sql-core}

{\small
\rowcolors{3}{rowshade}{white}
\setlength{\tabcolsep}{4.2pt}
\begin{longtable}{@{}lccrrrrrrrrrrr@{}}
\caption{Core study on Llama-3.1-8B base: every seed and round of Table-Specialist (TS) and no-validation (NV).}\label{tab:gen-perseed_core_base}\\
\toprule
\rowcolor{headshade}Arm & Seed & Rd & \multicolumn{2}{c}{WikiSQL} & \multicolumn{2}{c}{WikiTQ} & \multicolumn{2}{c}{Spider} & \multicolumn{2}{c}{BIRD} & M8 & ZS4 & Kept\\\rowcolor{headshade} & & & z & f & z & f & z & f & z & f & & &\\
\midrule
\endfirsthead
\multicolumn{14}{l}{\textit{\tablename~\thetable{} (continued)}}\\
\toprule
\rowcolor{headshade}Arm & Seed & Rd & \multicolumn{2}{c}{WikiSQL} & \multicolumn{2}{c}{WikiTQ} & \multicolumn{2}{c}{Spider} & \multicolumn{2}{c}{BIRD} & M8 & ZS4 & Kept\\\rowcolor{headshade} & & & z & f & z & f & z & f & z & f & & &\\
\midrule
\endhead
\midrule
\multicolumn{14}{r}{\textit{continued on next page}}\\
\endfoot
\bottomrule
\endlastfoot
TS & 1 & 1 & 0.6800 & 0.6490 & 0.2920 & 0.3120 & 0.7154 & 0.7037 & 0.3664 & 0.3664 & 0.5106 & 0.5134 & 426\\
TS & 1 & 2 & 0.7560 & 0.7500 & 0.2970 & 0.3250 & 0.7287 & 0.7237 & 0.3774 & 0.3609 & 0.5398 & 0.5398 & 608\\
TS & 1 & 3 & 0.7900 & 0.7720 & 0.3540 & 0.3380 & 0.7496 & 0.7337 & 0.4050 & 0.3857 & 0.5660 & 0.5746 & 672\\
TS & 2 & 1 & 0.7040 & 0.7330 & 0.2690 & 0.3110 & 0.6928 & 0.6953 & 0.3471 & 0.3747 & 0.5159 & 0.5032 & 388\\
TS & 2 & 2 & 0.7510 & 0.7490 & 0.2700 & 0.3000 & 0.6803 & 0.6811 & 0.3196 & 0.3361 & 0.5109 & 0.5052 & 655\\
TS & 2 & 3 & 0.7720 & 0.7640 & 0.2860 & 0.3080 & 0.6753 & 0.6711 & 0.3526 & 0.3664 & 0.5244 & 0.5215 & 777\\
TS & 3 & 1 & 0.7850 & 0.7800 & 0.3420 & 0.3420 & 0.7003 & 0.6861 & 0.3967 & 0.3939 & 0.5533 & 0.5560 & 423\\
TS & 3 & 2 & 0.7310 & 0.7240 & 0.3240 & 0.3310 & 0.7062 & 0.6953 & 0.3857 & 0.3829 & 0.5350 & 0.5367 & 618\\
TS & 3 & 3 & 0.7750 & 0.7690 & 0.3250 & 0.3440 & 0.6845 & 0.6594 & 0.3829 & 0.3802 & 0.5400 & 0.5419 & 729\\
\midrule
NV & 1 & 1 & 0.7100 & 0.6800 & 0.3510 & 0.3330 & 0.7604 & 0.7262 & 0.3912 & 0.3829 & 0.5418 & 0.5532 & 2797\\
NV & 1 & 2 & 0.6760 & 0.6390 & 0.3500 & 0.3250 & 0.7237 & 0.6953 & 0.3939 & 0.3857 & 0.5236 & 0.5359 & 2863\\
NV & 1 & 3 & 0.6920 & 0.6270 & 0.3300 & 0.3250 & 0.6820 & 0.6644 & 0.3747 & 0.3581 & 0.5067 & 0.5197 & 2854\\
NV & 2 & 1 & 0.6660 & 0.6650 & 0.3530 & 0.3310 & 0.7462 & 0.7137 & 0.4132 & 0.3857 & 0.5342 & 0.5446 & 2788\\
NV & 2 & 2 & 0.6380 & 0.6290 & 0.3300 & 0.3040 & 0.7020 & 0.6669 & 0.3857 & 0.3691 & 0.5031 & 0.5139 & 2852\\
NV & 2 & 3 & 0.6750 & 0.6420 & 0.3080 & 0.2830 & 0.6628 & 0.6361 & 0.3636 & 0.3416 & 0.4890 & 0.5023 & 2832\\
NV & 3 & 1 & 0.7000 & 0.7070 & 0.3670 & 0.3370 & 0.7538 & 0.7254 & 0.3939 & 0.3691 & 0.5442 & 0.5537 & 2797\\
NV & 3 & 2 & 0.6660 & 0.6770 & 0.3390 & 0.3200 & 0.7245 & 0.7062 & 0.4022 & 0.3664 & 0.5252 & 0.5329 & 2819\\
NV & 3 & 3 & 0.6210 & 0.5820 & 0.3100 & 0.3040 & 0.6953 & 0.6661 & 0.3774 & 0.3416 & 0.4872 & 0.5009 & 2840\\
\end{longtable}
}
\vspace{-0.6\baselineskip}
{\footnotesize\color{slate}\textit{Note.} Repo-faithful execution accuracy; z = zero-shot, f = few-shot prompt. M8 = mean of the 8 cells, ZS4 = mean of the 4 zero-shot cells. Kept = training pairs (validated pairs; every pair with code for NV; kept plus recovered for repair-TS). Source: \fpath{results/baseline/train_<arm>_<profile>_seed<s>_round<r>_*.json} and \fpath{runs/.../summary.json}.\par}
\medskip


\Cref{tab:gen-round-stats} summarises every NL-to-SQL approach of this and the next chapter over seeds, per round:
mean, sample standard deviation, best and worst seed.

{\small
\rowcolors{3}{rowshade}{white}
\setlength{\tabcolsep}{4.2pt}
\begin{longtable}{@{}lccrrrrrrrr@{}}
\caption{NL-to-SQL summary statistics over seeds for every approach and round. SD is the sample standard deviation ($n-1$); it is omitted when only one seed exists. Vanilla is a single deterministic run.}\label{tab:gen-round-stats}\\
\toprule
\rowcolor{headshade}Approach & Rd & Seeds & \multicolumn{4}{c}{M8 over seeds} & \multicolumn{4}{c}{ZS4 over seeds}\\\rowcolor{headshade} & & & Mean & SD & Min & Max & Mean & SD & Min & Max\\
\midrule
\endfirsthead
\multicolumn{11}{l}{\textit{\tablename~\thetable{} (continued)}}\\
\toprule
\rowcolor{headshade}Approach & Rd & Seeds & \multicolumn{4}{c}{M8 over seeds} & \multicolumn{4}{c}{ZS4 over seeds}\\\rowcolor{headshade} & & & Mean & SD & Min & Max & Mean & SD & Min & Max\\
\midrule
\endhead
\midrule
\multicolumn{11}{r}{\textit{continued on next page}}\\
\endfoot
\bottomrule
\endlastfoot
Base vanilla & -- & 1 & 0.4788 & -- & 0.4788 & 0.4788 & 0.4083 & -- & 0.4083 & 0.4083\\
Base TS & 1 & 3 & 0.5266 & 0.0232 & 0.5106 & 0.5533 & 0.5242 & 0.0280 & 0.5032 & 0.5560\\
Base TS & 2 & 3 & 0.5286 & 0.0155 & 0.5109 & 0.5398 & 0.5272 & 0.0191 & 0.5052 & 0.5398\\
Base TS & 3 & 3 & 0.5435 & 0.0210 & 0.5244 & 0.5660 & 0.5460 & 0.0268 & 0.5215 & 0.5746\\
Base NV & 1 & 3 & 0.5401 & 0.0052 & 0.5342 & 0.5442 & 0.5505 & 0.0051 & 0.5446 & 0.5537\\
Base NV & 2 & 3 & 0.5173 & 0.0123 & 0.5031 & 0.5252 & 0.5276 & 0.0119 & 0.5139 & 0.5359\\
Base NV & 3 & 3 & 0.4943 & 0.0108 & 0.4872 & 0.5067 & 0.5077 & 0.0104 & 0.5009 & 0.5197\\
Base repair-TS & 1 & 3 & 0.5509 & 0.0106 & 0.5388 & 0.5584 & 0.5530 & 0.0119 & 0.5393 & 0.5610\\
Base repair-TS & 2 & 3 & 0.5713 & 0.0169 & 0.5543 & 0.5881 & 0.5780 & 0.0192 & 0.5589 & 0.5972\\
Base repair-TS & 3 & 3 & 0.5856 & 0.0051 & 0.5801 & 0.5901 & 0.5904 & 0.0037 & 0.5867 & 0.5941\\
Base bestof4 & 1 & 3 & 0.5330 & 0.0141 & 0.5171 & 0.5441 & 0.5308 & 0.0184 & 0.5102 & 0.5455\\
Base DPO & 1 & 3 & 0.5057 & 0.0209 & 0.4859 & 0.5276 & 0.5051 & 0.0275 & 0.4740 & 0.5263\\
Base same-data control & 1 & 3 & 0.5385 & 0.0067 & 0.5320 & 0.5454 & 0.5359 & 0.0061 & 0.5300 & 0.5421\\
Instruct vanilla & -- & 1 & 0.5546 & -- & 0.5546 & 0.5546 & 0.5001 & -- & 0.5001 & 0.5001\\
Instruct TS & 1 & 3 & 0.6152 & 0.0106 & 0.6070 & 0.6272 & 0.6165 & 0.0100 & 0.6096 & 0.6280\\
Instruct TS & 2 & 3 & 0.6127 & 0.0172 & 0.5960 & 0.6304 & 0.6176 & 0.0166 & 0.6014 & 0.6346\\
Instruct TS & 3 & 3 & 0.6069 & 0.0131 & 0.5920 & 0.6169 & 0.6113 & 0.0135 & 0.5961 & 0.6220\\
Instruct NV & 1 & 1 & 0.6052 & -- & 0.6052 & 0.6052 & 0.6162 & -- & 0.6162 & 0.6162\\
Instruct NV & 2 & 1 & 0.5954 & -- & 0.5954 & 0.5954 & 0.6094 & -- & 0.6094 & 0.6094\\
Instruct NV & 3 & 1 & 0.5714 & -- & 0.5714 & 0.5714 & 0.5852 & -- & 0.5852 & 0.5852\\
Instruct repair-TS & 1 & 3 & 0.6117 & 0.0140 & 0.5955 & 0.6203 & 0.6156 & 0.0187 & 0.5941 & 0.6275\\
Instruct repair-TS & 2 & 3 & 0.6180 & 0.0026 & 0.6155 & 0.6208 & 0.6241 & 0.0047 & 0.6186 & 0.6269\\
Instruct repair-TS & 3 & 3 & 0.6054 & 0.0075 & 0.5972 & 0.6118 & 0.6122 & 0.0107 & 0.6009 & 0.6223\\
Instruct same-data control & 1 & 3 & 0.6231 & 0.0095 & 0.6166 & 0.6340 & 0.6269 & 0.0102 & 0.6184 & 0.6382\\
\end{longtable}
}
\vspace{-0.6\baselineskip}
{\footnotesize\color{slate}\textit{Note.} Computed by \fpath{report/tools/build_tables.py} from the per-cell result files; these are recalculated statistics, not values copied from the project log.\par}
\medskip


\begin{figure}[htbp]
\centering
\begin{tikzpicture}
\begin{groupplot}[group style={group size=2 by 1,horizontal sep=1.6cm},
  width=0.48\textwidth,height=6.2cm,xmin=0.7,xmax=3.3,xtick={1,2,3},xlabel={Round},
  ylabel={\Mviii{} (mean of 8 cells)},ymajorgrids,grid style={headshade},
  legend style={font=\scriptsize,at={(0.03,0.97)},anchor=north west,draw=none,fill=white,fill opacity=0.85},
  tick label style={font=\scriptsize},label style={font=\small},title style={font=\small\bfseries}]
\nextgroupplot[title={Llama-3.1-8B base},ymin=0.46,ymax=0.61]
  \addplot[navy,thick,mark=*,error bars/.cd,y dir=both,y explicit] coordinates {\plotBaseTS};
  \addplot[coral,thick,mark=square*,error bars/.cd,y dir=both,y explicit] coordinates {\plotBaseNV};
  \addplot[sea,thick,mark=triangle*,error bars/.cd,y dir=both,y explicit] coordinates {\plotBaseRTS};
  \addplot[slate,dashed,domain=0.7:3.3] {\valBaseVanilla};
  \legend{TS,NV,repair-TS,vanilla}
\nextgroupplot[title={Llama-3.1-8B-Instruct},ymin=0.54,ymax=0.66,legend style={at={(0.03,0.03)},anchor=south west}]
  \addplot[navy,thick,mark=*,error bars/.cd,y dir=both,y explicit] coordinates {\plotInstTS};
  \addplot[coral,thick,mark=square*] coordinates {\plotInstNV};
  \addplot[sea,thick,mark=triangle*,error bars/.cd,y dir=both,y explicit] coordinates {\plotInstRTS};
  \addplot[slate,dashed,domain=0.7:3.3] {\valInstVanilla};
  \legend{TS,NV (seed 1),repair-TS,vanilla}
\end{groupplot}
\end{tikzpicture}
\caption{NL-to-SQL \Mviii{} by round: mean over seeds 1--3, with error bars of one sample standard deviation (NV on
Instruct is one seed, so it has none). Plotted from the per-cell result files.}
\label{fig:rounds-sql}
\end{figure}

\begin{observation}
At round~3 on base (\cref{tab:gen-perseed_core_base,tab:gen-round-stats}):
\begin{itemize}
\item TS has \Mviii{} $0.5435\pm0.0210$ and NV $0.4943\pm0.0108$; vanilla is 0.4788.
\item TS over NV: $+0.0492$ [$+0.0426$, $+0.0557$], $p<0.0001$, positive on all three seeds ($+0.0593$, $+0.0354$,
$+0.0528$).
\item TS over vanilla: $+0.0647$ [$+0.0554$, $+0.0738$], $p<0.0001$, all seeds positive.
\end{itemize}
The orderings change across rounds (\cref{fig:rounds-sql}):
\begin{itemize}
\item In round~1, NV is higher than TS (0.5401 against 0.5266), and higher in 17 of the 24 (seed, cell) comparisons.
\item By round~3, TS is higher in 20 of 24.
\item NV falls every round (0.5401, 0.5173, 0.4943). TS rises (0.5266, 0.5286, 0.5435).
\end{itemize}
\end{observation}

The TS gain over vanilla is uneven across cells (\cref{tab:gen-significance}, ``cells''):
\begin{itemize}
\item Zero-shot gains: $+0.453$ on WikiSQL, $+0.103$ on WikiTQ, $+0.017$ on BIRD and $-0.021$ on Spider.
\item Few-shot cells change little or fall: $+0.083$ on WikiSQL, $-0.014$ on WikiTQ, $-0.081$ on Spider and $-0.022$
on BIRD.
\end{itemize}
Fine-tuning on zero-shot-format examples mostly teaches the zero-shot prompt's output format. Base vanilla already does
well with few-shot examples.

\begin{interpretation}
Two effects explain the round-1 advantage of NV and its later decline.
\begin{itemize}
\item \textbf{Volume.} NV trains on about 6.6 times as many pairs as TS in round~1 (2,797 against 426 for seed~1),
which helps when the model starts weak.
\item \textbf{Compounding errors.} From round~2, NV's generator is a model trained on unvalidated data, and its own
errors compound. TS's validated data keeps each round's generator from drifting.
\end{itemize}
The paper reports this ablation only as a figure, and does not show the round-1 crossover.
\end{interpretation}

\section{Validation outcomes}\label{sec:sql-validation}

\begin{landscape}
{\footnotesize
\rowcolors{3}{rowshade}{white}
\setlength{\tabcolsep}{4.2pt}
\begin{longtable}{@{}llccrrrrrrrrrrrr@{}}
\caption{Algorithm~3 outcomes for every validated round (3,000 tables per round). Each question yields one SQL--Scala pair; a pair is kept only if both programs run and agree on all 10 row samples. Rejection reasons are mutually exclusive and are assigned in the order the validator checks them.}\label{tab:gen-validation}\\
\toprule
\rowcolor{headshade}Model & Arm & S & R & Questions & SQL code & \makecell{Scala\\code} & Kept & Kept \% & \makecell{Always\\empty} & Disagree & \makecell{Scala\\failed} & \makecell{SQL\\failed} & No code & Error & \makecell{Scala\\runs}\\
\midrule
\endfirsthead
\multicolumn{16}{l}{\textit{\tablename~\thetable{} (continued)}}\\
\toprule
\rowcolor{headshade}Model & Arm & S & R & Questions & SQL code & \makecell{Scala\\code} & Kept & Kept \% & \makecell{Always\\empty} & Disagree & \makecell{Scala\\failed} & \makecell{SQL\\failed} & No code & Error & \makecell{Scala\\runs}\\
\midrule
\endhead
\midrule
\multicolumn{16}{r}{\textit{continued on next page}}\\
\endfoot
\bottomrule
\endlastfoot
Base & TS & 1 & 1 & 2904 & 2844 & 2848 & 426 & 14.7 & 621 & 840 & 855 & 53 & 105 & 4 & 6172\\
Base & TS & 1 & 2 & 2918 & 2900 & 2900 & 608 & 20.8 & 608 & 952 & 658 & 59 & 32 & 1 & 7926\\
Base & TS & 1 & 3 & 2918 & 2874 & 2888 & 672 & 23.0 & 622 & 900 & 559 & 95 & 67 & 3 & 8614\\
Base & TS & 2 & 1 & 2888 & 2846 & 2820 & 388 & 13.4 & 660 & 820 & 857 & 61 & 100 & 2 & 5769\\
Base & TS & 2 & 2 & 2916 & 2880 & 2873 & 655 & 22.5 & 621 & 942 & 538 & 83 & 73 & 4 & 8284\\
Base & TS & 2 & 3 & 2911 & 2891 & 2875 & 777 & 26.7 & 585 & 954 & 502 & 38 & 53 & 2 & 9576\\
Base & TS & 3 & 1 & 2893 & 2848 & 2837 & 423 & 14.6 & 656 & 781 & 876 & 57 & 96 & 4 & 6088\\
Base & TS & 3 & 2 & 2890 & 2875 & 2857 & 618 & 21.4 & 620 & 941 & 565 & 98 & 46 & 2 & 7921\\
Base & TS & 3 & 3 & 2908 & 2894 & 2849 & 729 & 25.1 & 618 & 951 & 485 & 52 & 70 & 3 & 9037\\
Base & repair-TS & 1 & 2 & 2918 & 2898 & 2838 & 690 & 23.6 & 583 & 906 & 537 & 108 & 93 & 1 & 8604\\
Base & repair-TS & 1 & 3 & 2918 & 2882 & 2842 & 761 & 26.1 & 612 & 817 & 545 & 74 & 106 & 3 & 9316\\
Base & repair-TS & 2 & 2 & 2916 & 2902 & 2874 & 713 & 24.5 & 688 & 889 & 515 & 56 & 51 & 4 & 8808\\
Base & repair-TS & 2 & 3 & 2911 & 2899 & 2849 & 829 & 28.5 & 556 & 833 & 568 & 53 & 70 & 2 & 10001\\
Base & repair-TS & 3 & 2 & 2890 & 2857 & 2837 & 701 & 24.3 & 633 & 890 & 530 & 53 & 81 & 2 & 8714\\
Base & repair-TS & 3 & 3 & 2908 & 2879 & 2821 & 896 & 30.8 & 579 & 772 & 478 & 76 & 104 & 3 & 10509\\
Instruct & TS & 1 & 1 & 2978 & 2923 & 2932 & 319 & 10.7 & 1007 & 626 & 691 & 234 & 97 & 4 & 4710\\
Instruct & TS & 1 & 2 & 2976 & 2960 & 2932 & 499 & 16.8 & 1050 & 586 & 580 & 205 & 55 & 1 & 6377\\
Instruct & TS & 1 & 3 & 2974 & 2947 & 2935 & 522 & 17.6 & 993 & 601 & 591 & 203 & 61 & 3 & 6692\\
Instruct & TS & 2 & 1 & 2983 & 2917 & 2949 & 332 & 11.1 & 1028 & 605 & 702 & 216 & 98 & 2 & 4827\\
Instruct & TS & 2 & 2 & 2980 & 2945 & 2934 & 486 & 16.3 & 956 & 562 & 641 & 257 & 75 & 3 & 6310\\
Instruct & TS & 2 & 3 & 2973 & 2948 & 2909 & 529 & 17.8 & 938 & 625 & 537 & 256 & 86 & 2 & 6675\\
Instruct & TS & 3 & 1 & 2982 & 2913 & 2942 & 332 & 11.1 & 1034 & 587 & 719 & 204 & 102 & 4 & 4845\\
Instruct & TS & 3 & 2 & 2966 & 2927 & 2907 & 486 & 16.4 & 1009 & 581 & 572 & 227 & 89 & 2 & 6216\\
Instruct & TS & 3 & 3 & 2967 & 2943 & 2871 & 526 & 17.7 & 967 & 555 & 556 & 245 & 115 & 3 & 6596\\
Instruct & repair-TS & 1 & 2 & 2976 & 2955 & 2936 & 500 & 16.8 & 1048 & 600 & 569 & 203 & 55 & 1 & 6396\\
Instruct & repair-TS & 1 & 3 & 2974 & 2952 & 2919 & 507 & 17.0 & 1005 & 654 & 521 & 212 & 73 & 2 & 6550\\
Instruct & repair-TS & 2 & 2 & 2980 & 2949 & 2954 & 492 & 16.5 & 979 & 574 & 601 & 276 & 54 & 4 & 6350\\
Instruct & repair-TS & 2 & 3 & 2973 & 2951 & 2929 & 525 & 17.7 & 989 & 649 & 533 & 211 & 64 & 2 & 6628\\
Instruct & repair-TS & 3 & 2 & 2966 & 2935 & 2917 & 479 & 16.1 & 1009 & 619 & 519 & 268 & 70 & 2 & 6145\\
Instruct & repair-TS & 3 & 3 & 2967 & 2944 & 2897 & 531 & 17.9 & 997 & 621 & 532 & 196 & 87 & 3 & 6738\\
\end{longtable}
}
\vspace{-0.6\baselineskip}
{\footnotesize\color{slate}\textit{Note.} ``Error'' counts pairs stopped by an unexpected exception (all were SQLite integer overflows). Repair-TS round~1 reuses TS round~1's validation, so it has no row of its own. Sources: \fpath{runs/train_<arm>_<profile>_seed<s>/round<r>/{summary,result}.json}.\par}
\medskip
\end{landscape}


\begin{observation}
\begin{itemize}
\item \textbf{Base TS} keeps 13.4--26.7\% of questions per round, rising with rounds (seed~1: 426, 608, 672 pairs).
\item \textbf{Instruct TS} keeps fewer: 319--332 pairs in round~1 and 486--529 in rounds 2--3.
\item \textbf{Main reasons for rejection.} \texttt{scala\_failed} and \texttt{disagree} on base, and
\texttt{always\_empty} on Instruct (about 1,000 per round).
\end{itemize}
\end{observation}

\begin{hypothesis}
Instruct writes more questions whose SQL returns nothing on every half-sample (\texttt{always\_empty}). This may
reflect questions with highly selective filters, which return no rows once half the table is removed. We did not test
this.
\end{hypothesis}

The hand inspection of the base 300-table run (\cref{sec:fail-validator}) gives a rough measure of validator precision:
\begin{itemize}
\item 39 of 50 kept pairs fit their question, 5 were doubtful, and 6 did not.
\item In 17 of 50 rejected pairs, the SQL may have been usable.
\end{itemize}

\section{Instruct: matching the paper}\label{sec:sql-instruct}

{\small
\rowcolors{3}{rowshade}{white}
\setlength{\tabcolsep}{4.2pt}
\begin{longtable}{@{}lccrrrrrrrrrrr@{}}
\caption{Core study on Llama-3.1-8B-Instruct: TS on seeds 1--3 and NV on seed~1, every round.}\label{tab:gen-perseed_core_instruct}\\
\toprule
\rowcolor{headshade}Arm & Seed & Rd & \multicolumn{2}{c}{WikiSQL} & \multicolumn{2}{c}{WikiTQ} & \multicolumn{2}{c}{Spider} & \multicolumn{2}{c}{BIRD} & M8 & ZS4 & Kept\\\rowcolor{headshade} & & & z & f & z & f & z & f & z & f & & &\\
\midrule
\endfirsthead
\multicolumn{14}{l}{\textit{\tablename~\thetable{} (continued)}}\\
\toprule
\rowcolor{headshade}Arm & Seed & Rd & \multicolumn{2}{c}{WikiSQL} & \multicolumn{2}{c}{WikiTQ} & \multicolumn{2}{c}{Spider} & \multicolumn{2}{c}{BIRD} & M8 & ZS4 & Kept\\\rowcolor{headshade} & & & z & f & z & f & z & f & z & f & & &\\
\midrule
\endhead
\midrule
\multicolumn{14}{r}{\textit{continued on next page}}\\
\endfoot
\bottomrule
\endlastfoot
TS & 1 & 1 & 0.8010 & 0.8030 & 0.4220 & 0.4350 & 0.8180 & 0.8155 & 0.4711 & 0.4518 & 0.6272 & 0.6280 & 319\\
TS & 1 & 2 & 0.8140 & 0.7960 & 0.4450 & 0.4540 & 0.8222 & 0.8114 & 0.4573 & 0.4435 & 0.6304 & 0.6346 & 499\\
TS & 1 & 3 & 0.7790 & 0.7530 & 0.4260 & 0.4360 & 0.8139 & 0.8047 & 0.4435 & 0.4380 & 0.6118 & 0.6156 & 522\\
TS & 2 & 1 & 0.8080 & 0.7980 & 0.4110 & 0.4250 & 0.7980 & 0.7788 & 0.4215 & 0.4160 & 0.6070 & 0.6096 & 332\\
TS & 2 & 2 & 0.7940 & 0.7860 & 0.4220 & 0.4270 & 0.8022 & 0.7830 & 0.4490 & 0.4298 & 0.6116 & 0.6168 & 486\\
TS & 2 & 3 & 0.7560 & 0.7630 & 0.4220 & 0.4260 & 0.7988 & 0.7713 & 0.4077 & 0.3912 & 0.5920 & 0.5961 & 529\\
TS & 3 & 1 & 0.7640 & 0.7720 & 0.4020 & 0.4210 & 0.8130 & 0.8013 & 0.4683 & 0.4490 & 0.6113 & 0.6118 & 332\\
TS & 3 & 2 & 0.7640 & 0.7530 & 0.4130 & 0.4220 & 0.7988 & 0.7796 & 0.4298 & 0.4077 & 0.5960 & 0.6014 & 486\\
TS & 3 & 3 & 0.8180 & 0.7890 & 0.4220 & 0.4230 & 0.8047 & 0.7938 & 0.4435 & 0.4408 & 0.6169 & 0.6220 & 526\\
\midrule
NV & 1 & 1 & 0.7810 & 0.7590 & 0.4230 & 0.4240 & 0.8063 & 0.7863 & 0.4545 & 0.4077 & 0.6052 & 0.6162 & 2881\\
NV & 1 & 2 & 0.7860 & 0.7540 & 0.4270 & 0.4070 & 0.7922 & 0.7654 & 0.4325 & 0.3994 & 0.5954 & 0.6094 & 2899\\
NV & 1 & 3 & 0.7260 & 0.6930 & 0.4130 & 0.3920 & 0.7805 & 0.7487 & 0.4215 & 0.3967 & 0.5714 & 0.5852 & 2861\\
\end{longtable}
}
\vspace{-0.6\baselineskip}
{\footnotesize\color{slate}\textit{Note.} Repo-faithful execution accuracy; z = zero-shot, f = few-shot prompt. M8 = mean of the 8 cells, ZS4 = mean of the 4 zero-shot cells. Kept = training pairs (validated pairs; every pair with code for NV; kept plus recovered for repair-TS). Source: \fpath{results/baseline/train_<arm>_<profile>_seed<s>_round<r>_*.json} and \fpath{runs/.../summary.json}.\par}
\medskip


\begin{figure}[htbp]
\centering
\begin{tikzpicture}
\begin{axis}[ybar,bar width=10pt,width=0.86\textwidth,height=6.2cm,ymin=0,ymax=0.95,
  symbolic x coords={WikiSQL,WikiTQ,Spider,BIRD},xtick=data,ylabel={Execution accuracy (zero-shot)},
  ymajorgrids,grid style={headshade},enlarge x limits=0.18,
  legend style={font=\scriptsize,at={(0.5,1.02)},anchor=south,legend columns=4,draw=none},
  nodes near coords,every node near coord/.append style={font=\tiny,rotate=90,anchor=west},
  tick label style={font=\small}]
\addplot[fill=slate!50,draw=slate] coordinates {(WikiSQL,\valInstVanzwikisql) (WikiTQ,\valInstVanzwikitq) (Spider,\valInstVanzspider) (BIRD,\valInstVanzbird)};
\addplot[fill=amber!45,draw=amber] coordinates {(WikiSQL,0.525) (WikiTQ,0.300) (Spider,0.670) (BIRD,0.388)};
\addplot[fill=navy!75,draw=navy,error bars/.cd,y dir=both,y explicit] coordinates {(WikiSQL,\valInstTSzwikisql) +- (0,\valInstTSzsdwikisql) (WikiTQ,\valInstTSzwikitq) +- (0,\valInstTSzsdwikitq) (Spider,\valInstTSzspider) +- (0,\valInstTSzsdspider) (BIRD,\valInstTSzbird) +- (0,\valInstTSzsdbird)};
\addplot[fill=amber,draw=amber!70!black] coordinates {(WikiSQL,0.816) (WikiTQ,0.449) (Spider,0.690) (BIRD,0.438)};
\legend{Our Instruct vanilla,Paper vanilla,Our Instruct TS r3 (mean $\pm$ SD),Paper TS}
\end{axis}
\end{tikzpicture}
\caption{Instruct against the paper's Table~3, zero-shot, per test set. Our TS bars are the mean over seeds 1--3 at
round~3, with one-SD error bars.}
\label{fig:inst-vs-paper}
\end{figure}

\begin{observation}
At round~3, Instruct TS has \ZSiv{} $0.6113\pm0.0135$ against the paper's 0.5983 (\cref{fig:inst-vs-paper}). Per
test set, ours minus the paper's:
\begin{itemize}
\item WikiSQL $-0.032$, WikiTQ $-0.026$, BIRD $-0.006$: all within the $\pm0.05$ gate;
\item Spider $+0.116$.
\end{itemize}
Gains over vanilla are WikiSQL $+0.266$, WikiTQ $+0.134$, BIRD $+0.029$ and Spider $+0.015$. The paper's are
$+0.291$, $+0.149$, $+0.050$ and $+0.020$: the same order, and 0.005--0.025 smaller. TS round~3 over vanilla on
\Mviii{}: $+0.0523$ [$+0.0424$, $+0.0622$], $p<0.0001$, all seeds positive.
\end{observation}

\begin{observation}
Rounds do not help Instruct.
\begin{itemize}
\item \ZSiv{} by round: $0.6165\pm0.0100$, $0.6176\pm0.0166$, $0.6113\pm0.0135$.
\item Round~3 against round~1 on \Mviii{}: $-0.0083$ [$-0.0133$, $-0.0032$], $p=0.0012$.
\item Seeds 1 and 2 peak at round~2; seed~3 rises at round~3.
\item The best round per seed averages 0.6245 \ZSiv{}, still above the paper.
\end{itemize}
Instruct NV (seed~1) falls every round (0.6162, 0.6094, 0.5852 \ZSiv{}). TS beats it at round~3 by $+0.0403$ on
\Mviii{} ($p<0.0001$), but this is a single seed.
\end{observation}

\begin{interpretation}
The Instruct track reproduces the paper's NL-to-SQL Table~3 approximately: the same method and test protocol, with
substitute data. The base track reproduces the paper's \emph{direction} but not its level. Because the paper does not
state which round Table~3 reports, both the last-round mean (0.6113) and the best-round mean (0.6245) are given. Both
are above the paper's 0.5983, entirely because of Spider.
\end{interpretation}

\section{Training cost}

\Cref{tab:gen-training,tab:gen-timings} in \cref{app:supp} give every adapter's size, steps, loss, time and memory, and
every round's stage times.
\begin{itemize}
\item A complete, uninterrupted TS round on base took 4.5--6.8~h, of which validation was 2.1--3.5~h.
\item An uninterrupted NV round took 6.5--7.8~h: NV has no validation, but training on about 2,800 pairs per adapter is about 6.6 times more training than TS.
\item Peak reserved GPU memory in the NL-to-SQL SFT runs was 18.9--20.5~GB.
\item Rounds that were resumed after a crash record only the time after the restart, so their totals are lower.
\end{itemize}
